\chapter{Manifolds and smooth maps}
\label{chap:manifolds}

In this chapter we set up the basic objects of the library: manifolds modeled on locally convex
spaces and $\Ck{k}$ maps between them. The calculus is the Michal--Bastiani calculus of
\cite{LCA}; the references for the manifold theory are \cite[\S I.3]{Nee06} and
\cite[Ch.~3]{GN26}, and \cite{Lan99} for the Banach case. Throughout, $E$, $F$, $G$ denote real
Hausdorff locally convex spaces. We stress from the start that \emph{no completeness is assumed}:
nothing in this chapter needs it. Tangent bundles, vector fields and manifolds with rough boundary
are treated in the next chapters.

\section{Manifolds}

\begin{definition}[The Michal--Bastiani groupoid]
  \label{def:mb-groupoid}
  \lean{LocallyConvex.Manifolds.mbGroupoid}
  \leanok
  The groupoid $\mathcal{G}_{\mb}(E)$ consists of the homeomorphisms $\varphi\colon U \to V$
  between open subsets of $E$ such that $\varphi$ is smooth on $U$ and $\varphi^{-1}$ is smooth
  on $V$, in the sense of Michal--Bastiani.
\end{definition}

That this is a structure groupoid (closed under composition, inversion and restriction, and
local in nature) follows from the chain rule and the locality of smoothness.

\begin{definition}[Manifolds]
  \label{def:mb-manifold}
  \lean{LocallyConvex.Manifolds.MBManifold, LocallyConvex.Manifolds.BanachManifold, LocallyConvex.Manifolds.HilbertManifold}
  \leanok
  \uses{def:mb-groupoid}
  A \emph{manifold modeled on $E$} is a topological space $M$ with an atlas of charts with
  values in $E$ whose transition maps belong to $\mathcal{G}_{\mb}(E)$
  \cite[Definition I.3.1]{Nee06}. A \emph{Banach manifold} is a manifold modeled on a normed
  space which is complete; a \emph{Hilbert manifold} is one modeled on a complete inner product
  space.
\end{definition}

In Lean, a manifold is a Mathlib \texttt{ChartedSpace} $E$ $M$ with \texttt{HasGroupoid} for
$\mathcal{G}_{\mb}(E)$; Mathlib's charted spaces accept an arbitrary topological model space, so
no new infrastructure is needed. Every chart comes with a \emph{preferred chart}
$c_x = \texttt{chartAt}\ x$ at each point $x \in M$, and we use it systematically.

\begin{remark}[The model space need not be complete]
  \label{rem:man-complete}
  We follow the convention of Glöckner--Neeb \cite[\S5.7]{GN26}: the model space of a manifold
  is not assumed to be complete. The formal theory of this and the following chapters (charts,
  smooth maps, tangent bundle, vector fields, connections, submanifolds) does not use
  completeness, since the calculus of \cite{LCA} replaces integrals in the model space by
  Hahn--Banach arguments. Results that need completeness (integrals with values in $E$,
  ordinary differential equations, exponential maps, inverse function theorems) say so in their
  hypotheses. In particular, a manifold modeled on an incomplete normed space is a manifold in
  our sense, but it is not a Banach manifold. Without completeness things do go wrong:
  \cite[Example 5.6.8]{GN26} is a Lie group modeled on an incomplete normed algebra that has no
  exponential function.
\end{remark}

On normed spaces, Michal--Bastiani smoothness agrees with Fréchet smoothness, so for complete
normed $E$ a Banach manifold is the classical object of \cite{Lan99}.

\begin{proposition}[Products]
  \label{prop:man-product}
  \lean{LocallyConvex.Manifolds.prod_mem_mbGroupoid, LocallyConvex.Analysis.SmoothMBOn.prodMap,
    LocallyConvex.Manifolds.smoothMBM_fst, LocallyConvex.Manifolds.smoothMBM_snd, LocallyConvex.Manifolds.ContMDiffMBOn.prodMk}
  \leanok
  \uses{def:mb-manifold, def:contmdiff}
  If $M$ and $N$ are manifolds modeled on $E$ and $F$, then $M \times N$, with the product
  charts, is a manifold modeled on $E \times F$. The projections onto $M$ and $N$ are smooth, and
  a pair $(f, g)$ of $\Ck{k}$ maps $P \to M$, $P \to N$ is a $\Ck{k}$ map $P \to M \times N$.
\end{proposition}

\begin{proof}
  \leanok
  The transition maps of the product atlas are products $\varphi \times \psi$ of transition
  maps, and a product of smooth diffeomorphisms between open sets is again one, since
  $(\varphi \times \psi) = (\varphi \circ \pi_1, \psi \circ \pi_2)$ and pairs of smooth maps are
  smooth. In charts the projections are the linear maps $\pi_1, \pi_2$.
\end{proof}

In Lean the model space of the product is Mathlib's type synonym \texttt{ModelProd} $E$ $F$ of
$E \times F$, which avoids instance diamonds when $M$ or $N$ is itself a product.

\section{Smooth maps}

For maps $f \colon M \to N$, open $U \subseteq M$ and charts $e$ of $M$, $e'$ of $N$, we denote
by
\[
  D(f, U, e, e') = e(\operatorname{dom} e) \cap e^{-1}\bigl(U \cap f^{-1}(\operatorname{dom}
  e')\bigr) \subseteq E
\]
the domain of the local representative $e' \circ f \circ e^{-1}$. It is open when $f$ is
continuous on $U$ (\texttt{LocallyConvex.Manifolds.isOpen\_chartDomain}).

\begin{definition}[$\Ck{k}$ maps between manifolds]
  \label{def:contmdiff}
  \lean{LocallyConvex.Manifolds.chartDomain, LocallyConvex.Manifolds.ContMDiffMBOn, LocallyConvex.Manifolds.ContMDiffMB,
    LocallyConvex.Manifolds.SmoothMBMOn, LocallyConvex.Manifolds.SmoothMBM}
  \leanok
  \uses{def:mb-manifold}
  A map $f \colon M \to N$ is $\Ck{k}$ on an open set $U \subseteq M$ if it is continuous on $U$
  and, for every $x \in U$, the local representative $c'_{f(x)} \circ f \circ c_x^{-1}$ in the
  preferred charts is $\Ck{k}$ on $D(f, U, c_x, c'_{f(x)})$. It is \emph{smooth} on $U$ if it is
  $\Ck{k}$ on $U$ for every $k$.
\end{definition}

Only the preferred charts appear in the definition. The next result says that this is no
restriction.

\begin{theorem}[Independence of the charts]
  \label{thm:man-contmdiff-chart}
  \lean{LocallyConvex.Manifolds.smoothMBOn_transition, LocallyConvex.Manifolds.smoothMBOn_transition_of_mem_maximalAtlas,
    LocallyConvex.Manifolds.ContMDiffMBOn.contDiffMBOn_chart,
    LocallyConvex.Manifolds.ContMDiffMBOn.contDiffMBOn_chart_of_mem_maximalAtlas,
    LocallyConvex.Manifolds.contMDiffMB_of_chartAt}
  \leanok
  \uses{def:contmdiff}
  If $f$ is $\Ck{k}$ on $U$, then its local representatives $e' \circ f \circ e^{-1}$ in arbitrary
  charts of the atlases (or of the maximal atlases) are $\Ck{k}$ on $D(f, U, e, e')$.
  Conversely, a continuous map whose local representatives in the preferred charts are
  $\Ck{k}$ near the center of each chart is $\Ck{k}$.
\end{theorem}

\begin{proof}
  \leanok
  The transition maps are smooth by definition of the atlas. Near a point $y = e(p)$ we write
  $e' f e^{-1} = (e' c'^{-1}) \circ (c' f c^{-1}) \circ (c e^{-1})$ with $c = c_p$,
  $c' = c'_{f(p)}$; the three factors are $\Ck{k}$ and the chain rule applies. Since $\Ck{k}$ is a
  local property, this suffices.
\end{proof}

\begin{theorem}[Composition]
  \label{thm:man-comp}
  \lean{LocallyConvex.Manifolds.ContMDiffMBOn.comp, LocallyConvex.Manifolds.SmoothMBMOn.comp, LocallyConvex.Manifolds.SmoothMBM.comp,
    LocallyConvex.Manifolds.smoothMBMOn_id, LocallyConvex.Manifolds.smoothMBMOn_const}
  \leanok
  \uses{def:contmdiff, thm:man-contmdiff-chart}
  If $f$ is $\Ck{k}$ on $U$, $g$ is $\Ck{k}$ on $V$ and $f(U) \subseteq V$, then $g \circ f$ is
  $\Ck{k}$ on $U$. Identities and constant maps are smooth.
\end{theorem}

\begin{proof}
  \leanok
  In the preferred charts at $x$ and $g(f(x))$, near a point we insert the preferred chart $d$ at
  the intermediate point and write the local representative of $g \circ f$ as the composition of
  the local representatives of $g$ and $f$ in the charts involving $d$. These are $\Ck{k}$ by
  Theorem~\ref{thm:man-contmdiff-chart}, and we conclude by the chain rule.
\end{proof}

\begin{proposition}[Consistency with the calculus]
  \label{prop:man-consistency}
  \lean{LocallyConvex.Manifolds.contMDiffMBOn_iff_contDiffMBOn}
  \leanok
  \uses{def:contmdiff}
  For a map $f \colon E \to F$ between the model spaces themselves (with their identity charts)
  and $U \subseteq E$ open, $f$ is $\Ck{k}$ on $U$ as a map between manifolds if and only if it
  is $\Ck{k}$ on $U$ in the sense of Michal--Bastiani.
\end{proposition}

\begin{proof}
  \leanok
  The preferred charts are identities, so $D(f, U, \id, \id) = U$ and the local representative
  is $f$ itself.
\end{proof}

\begin{remark}[Universes]
  \label{rem:man-universes}
  The model spaces $E$, $F$, $G$ live in a common universe, because the Michal--Bastiani
  derivative of a map defined on $E$ is a map defined on $E \times E$. The manifolds themselves
  are unrestricted. A consequence that will reappear below is that real-valued functions on
  $M$ are not treated as maps between manifolds (the model space $\R$ need not be in the
  universe of $E$), but through their local representatives.
\end{remark}
